% !TeX root = ../fundamentals.tex
% Section 2.2 -- Representations for mobility (the chapter's spine)
% Draft 1 (2026-07-21). Obeys WRITING_LAW.md + GLOSSARY.md. Citation/number ledger at the section end.

\section{Representations for mobility}
\label{sec:fund:repr}

How a place is represented determines what a model can learn from it. This section
follows one line of development, from symbolic identifiers to the check-in-level
representation this dissertation adopts, because the central argument of the work is
that the representation, more than the architecture, decides whether multi-task
learning helps.

The starting point is the one-hot identifier, a sparse vector that marks a place or a
category by position and encodes no relationship between them. Distributed
representations replace it. The skip-gram model learns dense vectors in which
proximity reflects co-occurrence, so that related symbols sit close together in a
continuous space \cite{mikolov2013word2vec}. Carrying that idea to graphs, DeepWalk
samples random walks and treats them as sentences, learning node embeddings from
local neighborhoods \cite{perozzi2014deepwalk}, and node2vec generalizes the walk so
that a single objective can favor either homophily or structural role
\cite{grover2016node2vec}. Graph neural networks then made the aggregation explicit.
The graph convolutional network combines a node's features with those of its
neighbors through a localized spectral rule \cite{kipf2017gcn}; graph attention
networks weight each neighbor by learned attention rather than by fixed degree
\cite{velivckovic2017graph}; and GraphSAGE learns aggregator functions that
generalize to nodes unseen during training, making the embedding inductive
\cite{hamilton2017graphsage}.

These methods still need a training signal. For places, labels are scarce, so the
representation is learned without them, by an information objective. Mutual
information between two high-dimensional variables can be estimated and maximized by
gradient descent \cite{belghazi2018mine}, and Deep InfoMax turns this into a
representation learner by maximizing the mutual information between a local feature
and a global summary of the input \cite{hjelm2019dim}. Deep Graph Infomax (DGI)
applies the principle to graphs: it learns node representations, unsupervised, by
maximizing the mutual information between local patch representations and a global
graph summary \cite{velickovic2019deep}. In this project DGI produces the place
embeddings used as input to the first models. Hierarchical Graph Infomax (HGI) builds
directly on it. HGI extends the infomax objective across a hierarchy and is trained by
maximizing the mutual information among the POI, region, and city levels, yielding
region-aware place embeddings \cite{huang2023hgi}. HGI is the place-level baseline
representation the later chapters measure against, and it is the direct base of the
representation the dissertation contributes.

A place embedding, however it is trained, shares one property: it assigns each place a
single fixed vector. That vector is the same whether the place is visited on a weekday
morning or a Saturday night, by a commuter or a tourist. Yet a place serves different
functions in different contexts, and a representation that cannot express this cannot
distinguish two visits to the same place that mean different things. CTLE makes the
point concrete by learning location embeddings that are context- and time-aware, so
that the vector for a place depends on the visit \cite{lin2021ctle}. This is the
limitation the rest of the dissertation responds to: a per-place vector is static
across visits, and moving below the place level, to the check-in itself, is the way
to make the representation reflect the visit.

Before that step, several general encoders supply the context a static place vector
omits. They are worth naming at their own origin, because they were developed outside
mobility and are imported as components. Time2Vec is a model-agnostic encoding of
time that combines a linear term with periodic terms, so that hour-of-day and
day-of-week enter a model as continuous features \cite{kazemi2019time2vec}. SIREN
shows that networks with sinusoidal activations represent continuous signals and their
derivatives well, which is the general basis for periodic coordinate encoders
\cite{sitzmann2020implicit}. Space2Vec encodes absolute position and spatial context
at multiple scales \cite{mai2020multiscalerepresentationlearningspatial}, and
Sphere2Vec encodes points on a sphere so that the encoding respects spherical
geometry rather than a flat approximation
\cite{mai2023sphere2vecgeneralpurposelocationrepresentation}. A related encoder uses
spherical harmonics together with sinusoidal representation networks, again to honor
the spherical domain that flat sine-and-cosine features distort
\cite{rußwurm2024geographiclocationencodingspherical}. Feature-wise linear modulation
(FiLM) is a general conditioning mechanism that scales and shifts a layer's features
according to a side input \cite{perez2018film}. In the models of Chapters~3 and~4 it
conditions the shared network on a per-task identity, so that one trunk can serve both
tasks. The context these encoders carry enters the model separately, and how it enters
is the turning point of this dissertation. Chapter~3 feeds a single place-level
embedding and adds only the per-task conditioning. Chapter~4 keeps that model but
decomposes the input, replacing the single embedding with separate spatial, temporal,
and categorical components combined at the input, and it is this change to the
representation, rather than any change to the architecture, that moves the results.

The check-in-level representation, Check2HGI, completes the line. It extends the
graph-infomax hierarchy with a fourth level below the place, the check-in, and is
trained without task labels in the same infomax spirit, so that each visit, rather
than each place, carries its own vector. Where a place embedding answers ``what is
this place,'' a check-in-level representation answers ``what is this visit,'' and
Chapter~5 develops this representation and the joint model built on it.
Table~\ref{tab:fund:lineage} traces the full progression, from DGI through HGI to the
check-in level and the joint model that uses it.

% ---------------------------------------------------------------------------
% CITATION / NUMBER LEDGER -- 2.2 (the spine; DGI -> HGI -> check-in level explicit)
%  mikolov2013word2vec   skip-gram distributed embeddings (one-hot -> dense) [new bib; ERRATA: CBIC church2017word2vec -> this]
%  perozzi2014deepwalk   DeepWalk: skip-gram on graph random walks [new bib]
%  grover2016node2vec    node2vec: biased-walk neighborhood objective
%  kipf2017gcn           GCN: localized spectral neighbor aggregation [new bib]
%  velivckovic2017graph  GAT: learned attention over neighbors [ERRATA: cite ICLR 2018 version]
%  hamilton2017graphsage GraphSAGE: inductive aggregators [new bib]
%  belghazi2018mine      MINE: MI estimable/maximizable by gradient descent [new bib]
%  hjelm2019dim          Deep InfoMax: local-global MI representation [new bib]
%  velickovic2019deep    DGI: node reps by local-patch/global-summary MI (base of lineage) [ERRATA: consolidate DGI triple-key]
%  feng2017poi2vec       POI2Vec: geographical influence in POI embedding (early mobility place-embedding)
%  huang2023hgi          HGI: MI across POI-region-city hierarchy (VERIFIED firsthand, PDF). Direct base.
%  kazemi2019time2vec    Time2Vec: linear+periodic time encoding (general, cited at origin) [new bib]
%  sitzmann2020implicit  SIREN: sinusoidal implicit neural reps (general)
%  mai2020...spatial     Space2Vec: multi-scale position encoder (general)
%  mai2023sphere2vec...  Sphere2Vec: spherical location encoder (general)
%  rußwurm2024...spherical spherical harmonics + SIREN location encoder (general; already in CoUrb bib -- existing key)
%  perez2018film         FiLM: feature-wise linear modulation. FIX (domain F3): in Ch.3/4 FiLM conditions on
%                        TASK IDENTITY (per-task gamma/beta); context enters by concatenation, NOT via FiLM.
%  lin2021ctle           CTLE: context/time-aware per-visit embedding = the static->contextual hinge
%  ENCODER ATTRIBUTION (domain F1 BLOCKER, FIXED): MTLnet (Ch.3) uses ONLY the place embedding + per-task FiLM;
%                        the decomposed spatial/temporal/categorical encoders are CoUrb's (Ch.4) contribution.
%                        Space2Vec (mai2020...) is named as background but adopted in NEITHER chapter.
%  Check2HGI / check-in-level representation: the MobiWac work (submitted, under review). NO bib cite here --
%                        silva2025mtlnet is the CBIC/MTLnet paper, NOT Check2HGI (fact-gate F1 BLOCKER, FIXED).
%                        Deferred to \ref{ch:mobiwac} + the lineage table, as the lineage table already does.
%  NUMBERS: none quoted in 2.2 (no dataset statistics). "fourth level", "seven"/counts not used here.
%  CROSS-REF: Table~\ref{tab:fund:lineage} = the model-lineage table.
% ---------------------------------------------------------------------------
